\section{Additional results}\label{app:more}

\begin{table}[h]\centering\small
\caption{Stress tests outside Assumption~\ref{as:unconf} and with misspecified
propensities (log-normal noise with $\sigma=0.5$). Relative $|$bias$|$ on unexposed
candidates.}\label{tab:stress}
\begin{adjustbox}{max width=\textwidth}
\begin{tabular}{lcccc}
\toprule
Scenario ($K=3$, rel.\ $|$bias$|$ on unexposed candidates) & IPS & IPS + WC & DR & \textbf{DRUP}\\
\midrule
independent exposures (Assumption~\ref{as:unconf}) & 0.029 & 0.029 & 1.137 & 0.019\\
fixed-size slates per user (Proposition~\ref{prop:dep}) & 0.219 & 0.219 & 1.119 & 0.024\\
misspecified propensities, exact imputation & 0.568 & 0.568 & 0.000 & 0.000\\
misspecified propensities, noisy imputation & 0.568 & 0.568 & 1.916 & 0.158\\
independent exposures (Assumption~\ref{as:unconf}), $K=5$ & 4.667 & 0.059 & 5.053 & 0.040\\
fixed-size slates per user (Proposition~\ref{prop:dep}), $K=5$ & 2.392 & 0.416 & 4.633 & 0.032\\
misspecified propensities, noisy imputation, $K=5$ & 19.266 & 1.162 & 16.027 & 0.204\\
\bottomrule
\end{tabular}

\end{adjustbox}
\end{table}

\begin{table}[h]\centering\small
\caption{Nuisance protocol in the 40-user, 60-item model of Table~\ref{tab:protocol}
with true and with estimated (logistic exposure model, misspecified here)
propensities, nuisances from an independent log, cross-fitted, or fitted on the
same log without cross-fitting, and degrees from $\hat Y$ or from edge
estimates. Relative $|$bias$|$ of the three-hop term and of the full score with
data-dependent constants (3{,}000 replications).}\label{tab:protocolfull}
\begin{adjustbox}{max width=\textwidth}
\begin{tabular}{lllcccccc}
\toprule
 & & & \multicolumn{3}{c}{IPS edges} & \multicolumn{3}{c}{DR edges}\\
Propensities & Nuisances & degrees & uncorr. & DRUP & $s$, data $c$ & uncorr. & DRUP & $s$, data $c$\\
\midrule
true & independent log & $\hat Y$ & 2.15 & 0.10 & 0.05 & 1.90 & 0.09 & 0.35\\
 & independent log & $W$ & 5.04 & 0.25 & 0.10 & 3.59 & 0.18 & 0.35\\
 & cross-fitted & $\hat Y$ & 0.39 & 0.64 & 0.15 & 0.61 & 0.80 & 0.62\\
 & cross-fitted & $W$ & 3.12 & 0.81 & 0.11 & 4.11 & 0.73 & 0.49\\
 & same log & $\hat Y$ & 0.42 & 0.70 & 0.18 & 0.59 & 0.70 & 0.33\\
 & same log & $W$ & 0.94 & 0.97 & 0.49 & 0.80 & 0.86 & 0.37\\
\midrule
estimated & independent log & $\hat Y$ & 29.57 & 3.70 & 0.15 & 25.57 & 2.11 & 0.53\\
 & independent log & $W$ & 92.66 & 7.44 & 0.15 & 47.74 & 3.00 & 0.50\\
 & cross-fitted & $\hat Y$ & 1.29 & 0.48 & 0.11 & 3.41 & 1.23 & 0.92\\
 & cross-fitted & $W$ & 83.73 & 0.49 & 0.18 & 27.63 & 0.73 & 0.59\\
 & same log & $\hat Y$ & 0.60 & 0.82 & 0.21 & 0.64 & 0.72 & 0.33\\
 & same log & $W$ & 0.93 & 0.97 & 0.53 & 0.83 & 0.88 & 0.38\\
\bottomrule
\end{tabular}

\end{adjustbox}
\end{table}

\begin{table}[h]\centering\small
\caption{Tightness of Theorem~\ref{thm:var} in a 12-user, 15-item model (20{,}000
replications, all propensities $\ge\tau$), over all $(u,i)$.}\label{tab:bounds}
\begin{adjustbox}{max width=\textwidth}\begin{tabular}{lcc}
\toprule
 & $\tau=0.3$ & $\tau=0.1$\\
\midrule
variance bound (b) / Monte-Carlo variance, median & 1582 & 8422\\
variance bound (b) / Monte-Carlo variance, minimum & 471 & 2126\\
McDiarmid half-width (d) / empirical 95\% quantile, median & 131 & 1831\\
McDiarmid half-width (d) / empirical 95\% quantile, minimum & 73 & 685\\
\bottomrule
\end{tabular}
\end{adjustbox}
\end{table}

\begin{table}[h]\centering\small
\caption{Exposure-capped re-ranking for all caps. Cells: nDCG@$K$ (Gini@$K$).}\label{tab:rerankfull}
\begin{adjustbox}{max width=\textwidth}\begin{tabular}{lccccccc}
\toprule
Coat & none & $c=20$ & $c=10$ & $c=5$ & $c=3$ & $c=2$ & $c=1.5$\\
\midrule
logged graph & 0.502 (0.34) & 0.494 (0.34) & 0.501 (0.34) & 0.501 (0.34) & 0.493 (0.33) & 0.494 (0.32) & 0.487 (0.28)\\
IPS adjacency & 0.507 (0.35) & 0.497 (0.35) & 0.506 (0.35) & 0.503 (0.34) & 0.504 (0.36) & 0.493 (0.33) & 0.486 (0.29)\\
DR adjacency & 0.489 (0.60) & 0.489 (0.60) & 0.489 (0.60) & 0.489 (0.60) & 0.485 (0.58) & 0.475 (0.48) & 0.444 (0.37)\\
DRUP & 0.489 (0.60) & 0.489 (0.60) & 0.489 (0.60) & 0.489 (0.60) & 0.485 (0.58) & 0.476 (0.48) & 0.443 (0.37)\\
\bottomrule
\end{tabular}
\end{adjustbox}

\medskip
\begin{adjustbox}{max width=\textwidth}\begin{tabular}{lccccccc}
\toprule
KuaiRec & none & $c=20$ & $c=10$ & $c=5$ & $c=3$ & $c=2$ & $c=1.5$\\
\midrule
logged graph & 0.603 (0.99) & 0.300 (0.91) & 0.228 (0.84) & 0.173 (0.72) & 0.139 (0.60) & 0.115 (0.44) & 0.102 (0.31)\\
IPS adjacency & 0.612 (0.99) & 0.307 (0.94) & 0.234 (0.89) & 0.178 (0.79) & 0.146 (0.66) & 0.121 (0.49) & 0.107 (0.34)\\
DR adjacency & 0.639 (0.99) & 0.335 (0.95) & 0.257 (0.90) & 0.199 (0.80) & 0.165 (0.67) & 0.140 (0.50) & 0.125 (0.35)\\
DRUP & 0.638 (0.99) & 0.335 (0.95) & 0.256 (0.90) & 0.199 (0.80) & 0.164 (0.67) & 0.140 (0.50) & 0.125 (0.35)\\
\bottomrule
\end{tabular}
\end{adjustbox}

\medskip
\begin{adjustbox}{max width=\textwidth}\begin{tabular}{lccccccc}
\toprule
\textit{Yahoo!\,R3}, cap factor $c$ & none & 20 & 10 & 5 & 3 & 2 & 1.5\\
\midrule
Obs & 0.658 (0.28) & 0.659 (0.28) & 0.660 (0.28) & 0.658 (0.28) & 0.656 (0.28) & 0.658 (0.27) & 0.652 (0.23)\\
IPS & 0.660 (0.27) & 0.663 (0.27) & 0.661 (0.27) & 0.662 (0.27) & 0.661 (0.27) & 0.659 (0.27) & 0.654 (0.22)\\
DR & 0.660 (0.27) & 0.663 (0.27) & 0.661 (0.27) & 0.662 (0.27) & 0.661 (0.27) & 0.659 (0.27) & 0.654 (0.22)\\
DRUP & 0.660 (0.27) & 0.663 (0.27) & 0.661 (0.27) & 0.662 (0.27) & 0.661 (0.27) & 0.659 (0.27) & 0.654 (0.22)\\
\bottomrule
\end{tabular}
\end{adjustbox}
\end{table}

\begin{table}[h]\centering\small
\caption{Best nDCG over each operator's hyper-parameter grid under a cap on the
Gini coefficient of top-$K$ exposure.}\label{tab:frontier}
\begin{adjustbox}{max width=\textwidth}
\begin{tabular}{llccccc}
\toprule
Data & Method & Gini$\le$0.3 & $\le$0.4 & $\le$0.5 & $\le$0.6 & any\\
\midrule
Coat & Obs & -- & 0.4917 & 0.4917 & 0.4917 & 0.4917\\
Coat & IPS & -- & 0.4864 & 0.4864 & 0.4864 & 0.4864\\
Coat & DR & -- & 0.4090 & 0.4487 & 0.4733 & 0.4853\\
Coat & DRUP & -- & 0.3786 & 0.4416 & 0.4848 & 0.4902\\
\bottomrule
\end{tabular}

\end{adjustbox}
\end{table}

\begin{table}[h]\centering\small
\caption{Certificates under re-fitted nuisances: the nuisances are re-fitted on the
poisoned log (no cross-fitting) and the realised top-$K$ is checked against the
frozen-nuisance certificates of the clean system.}\label{tab:refit}
\begin{adjustbox}{max width=\textwidth}
\begin{tabular}{llcccc}
\toprule
Data & Operator & fake users $F$ & hit rate (frozen / re-fitted) & certified pairs & violated\\
\midrule
Coat & IPS & 1, 2, 5, 10 & 0.001/0.001, 0.018/0.001, 0.268/0.001, 0.536/0.001 & 217 & 0\\
Coat & DR & 1, 2, 5, 10 & 0.000/0.000, 0.000/0.000, 0.000/0.000, 0.004/0.000 & 5193 & 0\\
Coat & DRUP & 1, 2, 5, 10 & 0.000/0.000, 0.000/0.000, 0.000/0.000, 0.004/0.000 & 5196 & 0\\
\bottomrule
\end{tabular}

\end{adjustbox}
\end{table}

\begin{table}[h]\centering\small
\caption{Jointly private DRUP for all fixed denoising ranks ($\delta=10^{-5}$).}\label{tab:privfull}
\begin{adjustbox}{max width=\textwidth}
\begin{tabular}{llccccccc}
\toprule
Data & rank & $\epsilon=0.5$ & 1 & 2 & 4 & 8 & 16 & $\infty$\\
\midrule
Coat & 8 & 0.381 & 0.385 & 0.400 & 0.399 & 0.398 & 0.449 & 0.533\\
 & 32 & 0.382 & 0.385 & 0.403 & 0.385 & 0.401 & 0.438 & 0.548\\
 & 128 & 0.372 & 0.370 & 0.387 & 0.401 & 0.389 & 0.422 & 0.536\\
 & full & 0.376 & 0.389 & 0.378 & 0.393 & 0.394 & 0.417 & 0.532\\
 & \multicolumn{8}{l}{\footnotesize non-private DRUP (all nuisances) 0.555; popularity 0.527; 29 public users}\\
\midrule
Yahoo!\,R3 & 8 & 0.483 & 0.492 & 0.491 & 0.493 & 0.492 & 0.493 & 0.492\\
 & 32 & 0.481 & 0.490 & 0.491 & 0.493 & 0.492 & 0.493 & 0.492\\
 & 128 & 0.479 & 0.490 & 0.490 & 0.494 & 0.493 & 0.492 & 0.493\\
 & full & 0.473 & 0.488 & 0.491 & 0.494 & 0.494 & 0.493 & 0.492\\
 & \multicolumn{8}{l}{\footnotesize non-private DRUP (all nuisances) 0.652; popularity 0.525; 1540 public users}\\
\midrule
KuaiRec & 8 & 0.175 & 0.179 & 0.173 & 0.427 & 0.603 & 0.618 & 0.634\\
 & 32 & 0.174 & 0.175 & 0.154 & 0.412 & 0.598 & 0.618 & 0.634\\
 & 128 & 0.160 & 0.173 & 0.172 & 0.365 & 0.581 & 0.614 & 0.634\\
 & full & 0.182 & 0.202 & 0.223 & 0.280 & 0.377 & 0.492 & 0.634\\
 & \multicolumn{8}{l}{\footnotesize non-private DRUP (all nuisances) 0.637; popularity 0.577; 717 public users}\\
\bottomrule
\end{tabular}

\end{adjustbox}
\end{table}

\begin{table}[h]\centering\small
\caption{Semi-synthetic KuaiRec with propensities estimated by the logistic
exposure model (columns as in Table~\ref{tab:semisynth}; regret not computed).}\label{tab:semisynthest}
\begin{adjustbox}{max width=\textwidth}
\begin{tabular}{llllccccc}
\toprule
 & & & & & \multicolumn{2}{c}{utility error (\%)} & & \\
$p$ & Nuisances & $C$ & Estimator & RMSE $T$ & random & popular & regret & nDCG@20\\
\midrule
est. & independent log & fixed & IPS & 77.7 & +361.4$\pm$5.6 & +138.2$\pm$1.1 & -- & 0.057\\
 &  &  & IPS+WC & 12.2 & +278.9$\pm$2.0 & +107.8$\pm$0.9 & -- & 0.057\\
 &  &  & DR & 92.2 & +370.0$\pm$7.6 & +146.3$\pm$1.1 & -- & 0.029\\
 &  &  & DRUP & 21.2 & +277.0$\pm$2.6 & +107.2$\pm$0.9 & -- & 0.038\\
\midrule
est. & cross-fitted & fixed & IPS & 78.4 & +359.9$\pm$6.9 & +136.1$\pm$1.0 & -- & 0.056\\
 &  &  & IPS+WC & 12.0 & +274.7$\pm$2.1 & +105.8$\pm$0.7 & -- & 0.056\\
 &  &  & DR & 113.8 & +430.5$\pm$8.2 & +155.3$\pm$1.1 & -- & 0.021\\
 &  &  & DRUP & 27.2 & +307.1$\pm$2.5 & +114.9$\pm$0.8 & -- & 0.030\\
\midrule
est. & sample split & fixed & IPS & 89.5 & +536.5$\pm$6.4 & +315.1$\pm$2.0 & -- & 0.066\\
 &  &  & IPS+WC & 12.0 & +453.5$\pm$2.6 & +280.8$\pm$1.8 & -- & 0.071\\
 &  &  & DR & 152.9 & +544.2$\pm$8.7 & +322.3$\pm$2.4 & -- & 0.020\\
 &  &  & DRUP & 23.4 & +446.6$\pm$3.0 & +276.0$\pm$2.0 & -- & 0.029\\
\midrule
\multicolumn{9}{l}{Five hops (independent log, true $p$, fixed $C$): relative $|$bias$|$ of $T_5$, DR 3.46, DRUP 0.32.}\\
\bottomrule
\end{tabular}

\end{adjustbox}
\end{table}

\begin{table}[h]\centering\small
\caption{Semi-synthetic KuaiRec with a sparser log (density 3\%, 58\% of the
propensities below the clip $\tau=0.01$), true and estimated propensities
(columns as in Table~\ref{tab:semisynth}; regret not computed).}\label{tab:semisynthsparse}
\begin{adjustbox}{max width=\textwidth}
\begin{tabular}{llllccccc}
\toprule
 & & & & & \multicolumn{2}{c}{utility error (\%)} & & \\
$p$ & Nuisances & $C$ & Estimator & RMSE $T$ & random & popular & regret & nDCG@20\\
\midrule
true & independent log & fixed & IPS & 326.6 & +49.1$\pm$26.9 & +23.2$\pm$1.1 & 0.816 & 0.083\\
 &  &  & IPS+WC & 11.4 & -27.9$\pm$2.5 & -12.1$\pm$0.5 & 0.816 & 0.083\\
 &  &  & DR & 341.1 & +331.4$\pm$28.5 & +121.2$\pm$1.7 & 0.839 & 0.037\\
 &  &  & DRUP & 39.9 & +163.6$\pm$5.4 & +64.2$\pm$1.2 & 0.841 & 0.046\\
\midrule
true & independent log & $\hat Y$ & IPS & -- & -49.3$\pm$2.3 & -42.4$\pm$0.4 & -- & 0.086\\
 &  &  & IPS+WC & -- & -64.7$\pm$1.2 & -54.2$\pm$0.2 & -- & 0.086\\
 &  &  & DR & -- & -8.9$\pm$2.5 & -30.7$\pm$0.4 & -- & 0.159\\
 &  &  & DRUP & -- & -30.2$\pm$1.3 & -42.9$\pm$0.3 & -- & 0.144\\
\midrule
true & cross-fitted & fixed & IPS & 326.6 & +49.1$\pm$26.9 & +23.2$\pm$1.1 & 0.816 & 0.083\\
 &  &  & IPS+WC & 11.4 & -27.9$\pm$2.5 & -12.1$\pm$0.5 & 0.816 & 0.083\\
 &  &  & DR & 521.4 & +484.6$\pm$30.7 & +166.3$\pm$3.5 & 0.851 & 0.030\\
 &  &  & DRUP & 73.3 & +260.3$\pm$8.1 & +108.2$\pm$3.0 & 0.848 & 0.040\\
\midrule
true & cross-fitted & $\hat Y$ & IPS & -- & -75.2$\pm$0.8 & -57.9$\pm$0.1 & -- & 0.098\\
 &  &  & IPS+WC & -- & -78.3$\pm$0.7 & -61.4$\pm$0.1 & -- & 0.098\\
 &  &  & DR & -- & -43.4$\pm$1.2 & -46.6$\pm$0.1 & -- & 0.140\\
 &  &  & DRUP & -- & -51.4$\pm$0.9 & -52.1$\pm$0.1 & -- & 0.133\\
\midrule
true & sample split & fixed & IPS & 291.2 & +200.4$\pm$26.7 & +197.0$\pm$5.6 & 0.855 & 0.052\\
 &  &  & IPS+WC & 9.0 & +123.3$\pm$3.9 & +158.9$\pm$5.3 & 0.855 & 0.049\\
 &  &  & DR & 3058.9 & +1572.9$\pm$142.3 & +848.2$\pm$69.1 & 0.865 & 0.020\\
 &  &  & DRUP & 294.4 & +932.8$\pm$38.9 & +583.1$\pm$29.9 & 0.863 & 0.029\\
\midrule
true & sample split & $\hat Y$ & IPS & -- & -10.1$\pm$5.3 & -19.5$\pm$1.4 & -- & 0.156\\
 &  &  & IPS+WC & -- & -43.8$\pm$1.6 & -39.6$\pm$0.7 & -- & 0.141\\
 &  &  & DR & -- & +24.8$\pm$5.1 & -8.5$\pm$1.4 & -- & 0.121\\
 &  &  & DRUP & -- & -13.3$\pm$1.6 & -28.7$\pm$0.8 & -- & 0.132\\
\midrule
est. & independent log & fixed & IPS & 930.7 & +333.6$\pm$32.8 & +122.0$\pm$1.7 & -- & 0.073\\
 &  &  & IPS+WC & 39.3 & +121.1$\pm$5.0 & +51.8$\pm$1.0 & -- & 0.073\\
 &  &  & DR & 957.3 & +891.6$\pm$39.0 & +268.3$\pm$2.3 & -- & 0.033\\
 &  &  & DRUP & 95.0 & +516.4$\pm$10.6 & +171.7$\pm$1.8 & -- & 0.040\\
\midrule
est. & cross-fitted & fixed & IPS & 798.1 & +331.0$\pm$33.1 & +119.1$\pm$1.7 & -- & 0.073\\
 &  &  & IPS+WC & 37.4 & +112.7$\pm$4.7 & +48.4$\pm$0.9 & -- & 0.073\\
 &  &  & DR & 1074.9 & +1242.2$\pm$48.4 & +322.0$\pm$4.6 & -- & 0.026\\
 &  &  & DRUP & 148.0 & +703.3$\pm$18.2 & +222.2$\pm$4.1 & -- & 0.036\\
\midrule
est. & sample split & fixed & IPS & 940.7 & +753.3$\pm$31.7 & +514.8$\pm$10.0 & -- & 0.062\\
 &  &  & IPS+WC & 40.0 & +502.0$\pm$8.2 & +415.5$\pm$8.9 & -- & 0.064\\
 &  &  & DR & 14661.2 & +2972.1$\pm$290.9 & +1342.7$\pm$83.3 & -- & 0.019\\
 &  &  & DRUP & 1239.1 & +1836.6$\pm$68.0 & +964.5$\pm$35.5 & -- & 0.028\\
\bottomrule
\end{tabular}

\end{adjustbox}
\end{table}

\begin{table}[p]\centering\scriptsize
\caption{AUC of the candidate list on unbiased test data (mean$\pm$std over splits; ties count one half). Same methods, tuning and splits as Table~\ref{tab:acc}, which tuned on the metric named there.}\label{tab:auc}
\begin{adjustbox}{max width=\textwidth}
\begin{tabular}{llcccc}
\toprule
 & Method & Coat AUC & Yahoo!\,R3 AUC & KuaiRand AUC & KuaiRec AUC\\
\midrule
\multicolumn{6}{l}{\emph{Non-personalised}}\\
 & Popularity & 0.644 & 0.655 & 0.581 & --\\
 & Imputation only & 0.654 & 0.652 & 0.546 & --\\
\multicolumn{6}{l}{\emph{Trained, pointwise loss}}\\
 & MF & 0.575 & 0.640 & -- & \textbf{0.802}\\
 & IPS-MF & 0.570 & 0.643 & -- & 0.789\\
 & DR-MF & 0.620 & 0.645 & -- & \underline{0.801}\\
 & DR-JL & 0.584 & 0.656 & -- & 0.797\\
 & MRDR & 0.589 & 0.657 & -- & 0.793\\
 & iALS & \textbf{0.667} & 0.761 & -- & 0.727\\
 & LightGCN & 0.632 & 0.705 & 0.559 & 0.799\\
 & DR-LightGCN & 0.658 & 0.675 & -- & 0.781\\
\multicolumn{6}{l}{\emph{Trained, pairwise or sampled loss}}\\
 & MF (BPR) & 0.614 & 0.751 & -- & 0.782\\
 & PDA & 0.612 & 0.758 & -- & 0.781\\
 & MACR & 0.639 & 0.729 & -- & 0.762\\
 & LightGCN (BPR) & 0.657 & \underline{0.771} & \underline{0.614} & 0.785\\
 & r-AdjNorm & 0.658 & \textbf{0.775} & \textbf{0.617} & 0.778\\
 & NAVIP & 0.653 & 0.771 & -- & 0.796\\
 & SimGCL & 0.646 & 0.736 & 0.578 & 0.765\\
\multicolumn{6}{l}{\emph{Training-free, logged graph}}\\
 & linear LightGCN / r-AdjNorm & 0.648 & 0.759 & 0.599 & --\\
 & EASE & 0.614 & 0.745 & 0.569 & --\\
 & GF-CF & 0.638 & 0.756 & 0.603 & --\\
 & BSPM & 0.665 & 0.756 & 0.604 & --\\
\multicolumn{6}{l}{\emph{Training-free, debiased graph}}\\
 & IPS adjacency (NAVIP-style) & 0.641 & 0.756 & -- & --\\
 & IPS + walk correction & 0.641 & 0.756 & -- & --\\
 & EASE on DR graph & 0.663 & 0.702 & -- & --\\
 & GF-CF on DR graph & 0.664 & 0.739 & -- & --\\
 & BSPM on DR graph & \underline{0.666} & 0.737 & -- & --\\
 & DR adjacency & 0.652 & 0.756 & -- & --\\
 & DRUP & 0.653 & 0.756 & -- & --\\
 & DR adjacency, 5 hops & 0.655 & -- & -- & --\\
 & DRUP, 5 hops & 0.649 & -- & -- & --\\
\multicolumn{6}{l}{\emph{Training-free, sample split (Assumption 2 holds)}}\\
 & DR adjacency, sample split & 0.659 & 0.708 & -- & --\\
 & DRUP-split & 0.658 & 0.708 & -- & --\\
\bottomrule
\end{tabular}

\end{adjustbox}
\end{table}

\begin{table}[p]\centering\scriptsize
\caption{Recall at the cut-off of Table~\ref{tab:acc} (@5 on Coat and Yahoo!\,R3, @10 on KuaiRand, @20 on KuaiRec). On Yahoo!\,R3 the lists have ten candidates, so Recall@10 is 1 for every method and Recall@5 is reported.}\label{tab:recall}
\begin{adjustbox}{max width=\textwidth}
\begin{tabular}{llcccc}
\toprule
 & Method & Coat Recall@5 & Yahoo!\,R3 Recall@5 & KuaiRand Recall@10 & KuaiRec Recall@20\\
\midrule
\multicolumn{6}{l}{\emph{Non-personalised}}\\
 & Popularity & 0.644 & 0.679 & 0.509 & 0.166\\
 & Imputation only & 0.663 & 0.678 & 0.480 & 0.186\\
\multicolumn{6}{l}{\emph{Trained, pointwise loss}}\\
 & MF & 0.579 & 0.670 & 0.469 & 0.189\\
 & IPS-MF & 0.573 & 0.671 & -- & 0.192\\
 & DR-MF & 0.624 & 0.675 & 0.481 & 0.188\\
 & DR-JL & 0.593 & 0.688 & 0.481 & 0.194\\
 & MRDR & 0.592 & 0.689 & -- & 0.193\\
 & iALS & \textbf{0.681} & 0.790 & 0.531 & 0.183\\
 & LightGCN & 0.638 & 0.731 & 0.494 & 0.193\\
 & DR-LightGCN & 0.671 & 0.704 & -- & 0.190\\
\multicolumn{6}{l}{\emph{Trained, pairwise or sampled loss}}\\
 & MF (BPR) & 0.622 & 0.778 & 0.531 & 0.178\\
 & PDA & 0.623 & 0.787 & -- & 0.176\\
 & MACR & 0.645 & 0.750 & 0.513 & 0.179\\
 & LightGCN (BPR) & 0.659 & 0.799 & \underline{0.542} & 0.180\\
 & r-AdjNorm & 0.662 & \textbf{0.802} & \textbf{0.543} & 0.180\\
 & NAVIP & 0.661 & \underline{0.800} & 0.541 & 0.185\\
 & SimGCL & 0.656 & 0.755 & 0.508 & 0.162\\
\multicolumn{6}{l}{\emph{Training-free, logged graph}}\\
 & linear LightGCN / r-AdjNorm & 0.646 & 0.781 & 0.529 & 0.181\\
 & EASE & 0.620 & 0.767 & 0.509 & 0.154\\
 & GF-CF & 0.657 & 0.777 & 0.532 & 0.183\\
 & BSPM & 0.675 & 0.779 & 0.531 & 0.183\\
\multicolumn{6}{l}{\emph{Training-free, debiased graph}}\\
 & IPS adjacency (NAVIP-style) & 0.645 & 0.783 & 0.528 & 0.185\\
 & IPS + walk correction & 0.645 & 0.783 & 0.528 & 0.185\\
 & propagation over $\hat Y$ only (no residual) & 0.664 & 0.699 & -- & 0.192\\
 & EASE on DR graph & 0.668 & 0.720 & 0.487 & 0.180\\
 & GF-CF on DR graph & 0.672 & 0.760 & 0.526 & \textbf{0.194}\\
 & BSPM on DR graph & \underline{0.677} & 0.756 & 0.528 & 0.194\\
 & DR adjacency & 0.654 & 0.783 & 0.528 & 0.194\\
 & DRUP & 0.655 & 0.783 & 0.528 & \underline{0.194}\\
 & DR adjacency, 5 hops & 0.662 & -- & -- & 0.194\\
 & DRUP, 5 hops & 0.655 & -- & -- & 0.194\\
\multicolumn{6}{l}{\emph{Training-free, sample split (Assumption 2 holds)}}\\
 & DR adjacency, sample split & 0.659 & 0.745 & 0.521 & 0.193\\
 & DRUP-split & 0.658 & 0.744 & 0.521 & 0.193\\
\bottomrule
\end{tabular}

\end{adjustbox}
\end{table}

\begin{table}[t]\centering\small
\caption{Which nuisance breaks the correction under cross-fitting. Relative absolute bias of the three-hop term
(40 users, 60 items, true propensities, 1{,}500 replications) with the degree weights $C$ fixed a priori
(degrees of $Y$), taken from the imputation $\hat Y$, or from the edge estimates $W$; nuisances from an
independent log, cross-fitted, fitted in sample, or from the 20\% sample split. Compare IPS and DR with
and without the walk correction.}\label{tab:protocolfixed}
\begin{adjustbox}{max width=\textwidth}
\begin{tabular}{llcccc}
\toprule
Nuisances & Degree weights $C$ & IPS, uncorrected & IPS+WC & DR, uncorrected & DRUP\\
\midrule
independent log & fixed a priori & 1.10 & 0.05 & 0.81 & 0.04\\
 & from $\hat Y$ & 2.20 & 0.13 & 1.95 & 0.11\\
 & from $W$ & 5.31 & 0.33 & 3.58 & 0.22\\
\midrule
cross-fitted & fixed a priori & 1.10 & 0.05 & 0.89 & 0.05\\
 & from $\hat Y$ & 0.39 & 0.64 & 0.61 & 0.80\\
 & from $W$ & 3.22 & 0.82 & 3.95 & 0.74\\
\midrule
in-sample & fixed a priori & 1.10 & 0.05 & 0.51 & 0.23\\
 & from $\hat Y$ & 0.43 & 0.71 & 0.59 & 0.70\\
 & from $W$ & 0.94 & 0.97 & 0.80 & 0.87\\
\midrule
sample split 20\% & fixed a priori & 0.89 & 0.07 & 0.58 & 0.07\\
 & from $\hat Y$ & 1.10 & 0.21 & 0.92 & 0.21\\
\bottomrule
\end{tabular}

\end{adjustbox}
\end{table}

\begin{table}[t]\centering\scriptsize
\caption{Key comparisons under three tests. $\Delta$: mean per-user difference of the dataset metric (first method minus second).
$p_{\mathrm{Holm}}$: user-level paired $t$-test, Holm-adjusted over all comparisons of the dataset.
$p_{\mathrm{NB}}$: corrected resampled $t$-test over the split means~\cite{nadeau2003}, Holm-adjusted over the same family. $p_{\mathrm{TOST}}$: two
one-sided user-level tests of equivalence within $\pm0.005$ nDCG, with the margins $\pm0.002$ and $\pm0.01$ in brackets as sensitivity; a small value means that
the two methods are equivalent. ``Identical rankings'': the two operators select the same configuration on every split, so no test applies. The margin is a convention, not derived from the data. Sample-split rows use the first draw.}\label{tab:nbtost}
\begin{adjustbox}{max width=\textwidth}
\begin{tabular}{lcccc}
\toprule
Comparison ($\Delta$; $p_{\mathrm{Holm}}$ / $p_{\mathrm{NB}}$ / $p_{\mathrm{TOST}}$) & Coat & Yahoo!\,R3 & KuaiRand & KuaiRec\\
\midrule
DRUP vs DR adjacency & +0.0014; 1.000 / 0.917 / 0.034 & +0.0000; 1.000 / 1.000 / $<10^{-3}$ & +0.0000; 1.000 / 1.000 / $<10^{-3}$ & -0.0001; 1.000 / 0.840 / $<10^{-3}$\\
DRUP vs DR adjacency, 5 hops & -0.0081; 1.000 / 0.872 / 0.809 & -- & -- & -0.0003; 0.170 / 0.219 / $<10^{-3}$\\
DRUP-split vs DR-split & -0.0033; 1.000 / 0.903 / 0.295 & -0.0012; 1.000 / 0.866 / 0.001 & +0.0002; 1.000 / 0.971 / $<10^{-3}$ & -0.0001; 1.000 / 0.311 / $<10^{-3}$\\
DRUP vs GF-CF on DR graph & -0.0143; 0.106 / 0.681 / 0.984 & +0.0247; $<10^{-3}$ / 0.037 / 1.000 & +0.0025; 0.043 / 0.011 / $<10^{-3}$ & +0.0004; 1.000 / 0.774 / $<10^{-3}$\\
DRUP-split vs DRUP & +0.0087; 1.000 / 0.780 / 0.625 & -0.0659; $<10^{-3}$ / $<10^{-3}$ / 1.000 & -0.0092; $<10^{-3}$ / 0.001 / 0.996 & -0.0035; 0.004 / 0.058 / 0.051\\
DRUP vs DR-JL & +0.0658; $<10^{-3}$ / 0.274 / 1.000 & +0.1350; $<10^{-3}$ / $<10^{-3}$ / 1.000 & +0.0739; $<10^{-3}$ / $<10^{-3}$ / 1.000 & +0.0044; 0.003 / 0.142 / 0.292\\
DRUP vs LightGCN (pointwise) & +0.0073; 1.000 / 0.853 / 0.585 & +0.0616; $<10^{-3}$ / 0.012 / 1.000 & +0.0443; $<10^{-3}$ / $<10^{-3}$ / 1.000 & +0.0103; $<10^{-3}$ / $<10^{-3}$ / 1.000\\
DRUP vs iALS & -0.0372; 1.000 / 0.362 / 0.968 & -0.0065; 1.000 / 0.174 / 0.643 & -0.0031; 1.000 / 0.134 / 0.119 & +0.0249; $<10^{-3}$ / $<10^{-3}$ / 1.000\\
DRUP vs r-AdjNorm & -0.0290; 1.000 / 0.614 / 0.920 & -0.0187; $<10^{-3}$ / 0.032 / 0.999 & -0.0101; $<10^{-3}$ / 0.019 / 0.999 & +0.0296; $<10^{-3}$ / $<10^{-3}$ / 1.000\\
\bottomrule
\end{tabular}

\end{adjustbox}
\end{table}

\subsection{Does the correction matter on sparser logs?}\label{sec:sparse}

\begin{table}[t]\centering\small
\caption{KuaiRec with the log thinned to a fraction $q$ of its exposures
(nuisances re-fitted, additive imputation, $\tau\in\{0.05,0.01\}$, $\lambda\in\{0,1\}$,
selected on validation). nDCG@20 of the uncorrected and corrected operators, and
the lowest top-20 overlap and Spearman correlation between the two rankings at
matched configurations.}\label{tab:sparse}
\begin{adjustbox}{max width=\textwidth}
\begin{tabular}{cccccccc}
\toprule
thinning $q$ & density & DR & DRUP & DR 5 hops & DRUP 5 hops & top-20 overlap & Spearman\\
\midrule
1 & 13.38\% & 0.6302 & 0.6302 & 0.6306 & 0.6306 & 0.971 & 0.999\\
0.1 & 1.34\% & 0.6097 & 0.6102 & 0.6109 & 0.6111 & 0.935 & 0.999\\
0.03 & 0.40\% & 0.6073 & 0.6069 & 0.6085 & 0.6081 & 0.985 & 1.000\\
\bottomrule
\end{tabular}

\end{adjustbox}
\end{table}

Theorem~\ref{thm:lower} predicts a larger bias when propensities are small. We
therefore thinned the KuaiRec log to 10\% and 3\% of its exposures, which lowers the
median logged propensity from 0.48 to 0.045 and 0.013, re-fitted the nuisances
and re-selected every operator (Table~\ref{tab:sparse}). The corrected and the
uncorrected operators still rank almost identically: their top-20 lists overlap
by at least 93\%, the rank correlation is at least 0.999, and the accuracy
differences stay below 0.0005 nDCG@20 in either direction, at three and at five
hops. On these data the repeated-walk term is small compared with the score
differences that decide the ranking, and the value of the correction lies in
the properties of the scores, not in the rankings they induce.


\subsection{Clip and exposure invariance; exposure concentration}\label{sec:fairsupp}
\begin{table}[t]\centering\small
\caption{Clip $\tau$ against exposure invariance for DRUP under the same
intervention, with frozen nuisances ($\eta$: mean$\pm$std over draws).}\label{tab:tau}
\begin{adjustbox}{max width=\textwidth}
\begin{tabular}{llcccc}
\toprule
Data & operator & $\tau$ & nDCG & score elasticity $\eta$ & rank shift\\
\midrule
Coat & Obs & -- & 0.4917 & +1.034$\pm$0.043 & -0.0646\\
Coat & DRUP & 0.2 & 0.4803 & -0.020$\pm$0.005 & +0.0061\\
Coat & DRUP & 0.05 & 0.4839 & -0.069$\pm$0.018 & +0.0188\\
Coat & DRUP & 0.01 & 0.4188 & -0.094$\pm$0.060 & +0.0212\\
Coat & DRUP & 0.001 & 0.4107 & -0.044$\pm$0.057 & +0.0121\\
\midrule
KuaiRec & Obs & -- & 0.6026 & +0.997$\pm$0.003 & -0.2412\\
KuaiRec & DRUP & 0.1 & 0.6305 & -0.013$\pm$0.000 & +0.0035\\
KuaiRec & DRUP & 0.02 & 0.3393 & +0.001$\pm$0.004 & +0.0027\\
KuaiRec & DRUP & 0.001 & 0.1942 & -0.045$\pm$0.022 & +0.0030\\
\midrule
Yahoo!\,R3 & Obs & -- & 0.6590 & +1.000$\pm$0.011 & -0.1604\\
Yahoo!\,R3 & DRUP & 0.2 & 0.6616 & +0.954$\pm$0.012 & -0.1507\\
Yahoo!\,R3 & DRUP & 0.05 & 0.6551 & +0.814$\pm$0.025 & -0.1282\\
Yahoo!\,R3 & DRUP & 0.02 & 0.6385 & +0.652$\pm$0.022 & -0.1070\\
Yahoo!\,R3 & DRUP & 0.005 & 0.6248 & +0.254$\pm$0.006 & -0.0630\\
Yahoo!\,R3 & DRUP & 0.002 & 0.6195 & +0.060$\pm$0.011 & -0.0468\\
Yahoo!\,R3 & DRUP & 0.001 & 0.6205 & +0.012$\pm$0.035 & -0.0440\\
Yahoo!\,R3 & DRUP & 0.0005 & 0.6204 & +0.006$\pm$0.039 & -0.0435\\
Yahoo!\,R3 & DRUP & 0.0002 & 0.6204 & +0.006$\pm$0.039 & -0.0435\\
\bottomrule
\end{tabular}

\end{adjustbox}
\end{table}

\begin{table}[t]\centering\small
\caption{Exposure concentration and user groups. PRU: popularity--rank
correlation. ECB (KuaiRec, where item quality is known from the fully observed
matrix): partial correlation between an item's mean rank and its log exposure,
controlling for its true like rate. User gap: nDCG of group 0 minus group 1
(men minus women on Coat, inactive minus active users elsewhere) with a
bootstrap 95\% interval.}\label{tab:fair}
\begin{adjustbox}{max width=\textwidth}
\begin{tabular}{llcccccc}
\toprule
Data & Method & nDCG & ECB$\downarrow$ & PRU$\downarrow$ & Gini$\downarrow$ & user gap$\downarrow$ & rank shift under $do(p/2)$\\
\midrule
Coat & Obs & 0.4917 & 0.769 & 0.394 & 0.340 & 0.1123 & -0.0585$\pm$0.0073\\
Coat & IPS & 0.4864 & 0.765 & 0.392 & 0.337 & 0.1189 & -0.0594$\pm$0.0066\\
Coat & DR & 0.4770 & 0.744 & 0.763 & 0.655 & 0.0795 & -0.0027$\pm$0.0122\\
Coat & DRUP & 0.4803 & 0.739 & 0.758 & 0.655 & 0.0793 & +0.0006$\pm$0.0110\\
\bottomrule
\end{tabular}

\end{adjustbox}
\end{table}

Table~\ref{tab:fair} shows the distributional side. On KuaiRec, the
exposure-conditional bias drops from 0.83 (logged graph) and 0.80 (IPS) to 0.33
for DRUP and DR adjacency: given their true quality, items are ranked much less
by how often they were shown. At the same time DRUP's top-20 lists are
concentrated on 1.2\% of the candidate items, against 15\% for the logged graph,
because the estimated full-exposure scores are themselves concentrated on few items; on Coat its Gini
coefficient is 0.60 against 0.34. Exposure invariance and exposure concentration
are different properties, and DRUP improves the first at the expense of the
second. The user-group gaps do not differ between operators. They are clearly
non-zero on Coat (men 0.11 above women), on Yahoo!\,R3 (inactive users 0.046
below active ones) and on KuaiRand (inactive users 0.014 above active ones), under every
operator, and not on KuaiRec. On KuaiRand the lists of DRUP, IPS and the logged graph
are almost equally concentrated (Gini 0.43, coverage 0.98), as the scores of all
three are close to each other there.


\subsection{Stakeholders and regulation}\label{sec:stake}
\emph{Item providers.} Exposure invariance means that an item's expected score does not
depend on how often it was observed; it is not equal exposure, and it comes with more
concentrated lists (Table~\ref{tab:fair}) that a retrained system could entrench.
Deployments should report concentration next to invariance and use an explicit allocation
mechanism. \emph{Users.} The removal attributions describe the model with its nuisances
frozen and are not an account of data deletion. \emph{Platforms and auditors.} The clip
$\tau$ and the control-variate weight $\lambda$ trade invariance against accuracy; one
option is to select the most accurate configuration subject to a bound on the elasticity
measured in a thinning audit, whose power is limited. The audit is a first-party
instrument: it needs the log, the propensities, the deployed hyper-parameters and the
ability to re-run the pipeline. Article~40 of the Digital Services
Act~\citep{dsa2022} gives vetted researchers access to data of very large online platforms
and search engines; whether it extends to re-running a proprietary pipeline is not settled
by the text, and we do not claim that it does. The Act's risk-assessment and audit
obligations (Articles~34--37, Delegated Regulation (EU) 2024/436~\citep{eu2024audit})
concern compliance of very large providers, not properties of estimators, and our results
can at most inform a risk assessment.
